\chapter{Computable Numbers and Functions}\label{ch:foundations}

A number will be a computation, not an already available point of a
completed line.  The computation gives rational bounds, and a proof explains
why they agree and how they become arbitrarily close.  Different computations
of the same number remain worth keeping: an elementary construction can
explain a quantity which a different construction computes more quickly.

The finite arithmetic used below is ordinary integer and rational arithmetic.
Fractions are reduced by the Euclidean algorithm; finite sums, products,
polynomials, induction, and the triangle inequality require no limiting
objects.  A rational input is a numerator and a positive denominator, not an
oracle for a point on a continuum.

\section{Interval arithmetic and real numbers}

\begin{definition}[Rational interval]\label{def:qinterval}
A \emph{rational interval} is a pair
\[
 [a,b]\in\Q\times\Q,\qquad a\le b.
\]
It is a finite object, not the set of rational points between its endpoints.
Its width is $b-a$.  Two intervals $[a,b]$ and $[c,d]$ overlap when
$a\le d$ and $c\le b$.
\end{definition}

\begin{definition}[A computed number]\label{def:raw-real}
A number computation is a terminating algorithm
\[
 n\longmapsto X_n=[a_n,b_n]
\]
whose intervals satisfy
\[
 n\le m\ \Longrightarrow\ a_n\le a_m\le b_m\le b_n,
\]
and whose widths become smaller than every positive rational tolerance.
We also call such an algorithm a raw real.  No infinitely stored table of
rational endpoints is allowed.
\end{definition}

A request for precision $\varepsilon>0$ is answered by computing successive
stages until $b_n-a_n<\varepsilon$.  A proof that the widths shrink proves
that this search terminates.  An explicit rate is often useful, but need not
be part of the algorithm's input.

\begin{definition}[Equivalence]\label{def:raw-equiv}
Two valid computations are equivalent when their intervals overlap at every
stage.  We write $X\sim Y$, and use the usual equality sign in mathematical
identities between their values.
\end{definition}

Same-stage overlap implies overlap at any two stages, by passing to their
maximum.  Equivalence is transitive: if two outer computations had disjoint
intervals separated by a positive rational gap, a sufficiently narrow interval
of a middle computation could not overlap both.  This proves transitivity
using only finitely many rational comparisons.

Equality is a theorem about computations, not a proposed decision algorithm.
To show two values differ, disjoint finite enclosures suffice.  To show a
value is positive, it suffices to find an enclosure with positive lower
endpoint.  Such a witness is called a separation from zero.

\begin{definition}[Interval arithmetic]\label{def:raw-real-arithmetic}
Addition and negation are computed by
\[
 [a,b]+[c,d]=[a+c,b+d],\qquad -[a,b]=[-b,-a].
\]
For multiplication take the least rational interval containing
$ac,ad,bc,bd$.  If an interval $[a,b]$ avoids zero, its reciprocal enclosure
is $[1/b,1/a]$.
\end{definition}

These operations preserve inclusion.  If the factors have absolute bounds
$B_X,B_Y$, the product width is at most
\[
 B_X\operatorname{width}(Y)+B_Y\operatorname{width}(X).
\]
The bounds can be read from any initial enclosures.  If $|X|\ge c>0$ is
certified, reciprocal widths are at most $c^{-2}\operatorname{width}(X)$
once the intervals have that separation.  These estimates prove validity
and compatibility with equivalence.  In particular, division is available
when a nonzero denominator has been certified, not by deciding an arbitrary
computed value's equality with zero.

An inequality such as $|X-Y|\le r$ has an interval meaning: arbitrarily
accurate enclosures certify the bound $r+\eta$ for each rational $\eta>0$.
An assertion of containment in a closed rational interval can equivalently
be phrased as overlap with that interval at every stage.  None of this
requires choosing a point in an infinite intersection.

\begin{proposition}[Finite-prefix stabilization]\label{prop:raw-prefix-stabilization}
Suppose rational numbers $q_n$ and positive rational errors $e_n\to0$
are computed, and
\[
 |q_m-q_n|\le e_n\qquad(m\ge n).
\]
Then
\[
 X_n=\bigcap_{k=0}^n[q_k-e_k,q_k+e_k]
\]
is a valid number computation.
\end{proposition}
\begin{proof}
The finite intersection contains $q_n$.  It is nested, and its width is at
most $2e_n$.  All its endpoints are finite maxima and minima of rationals.
\end{proof}

More generally, a sequence of pairwise overlapping rational intervals with
shrinking widths can be stabilized in this way: pairwise overlap of a
\emph{finite} family of intervals implies a nonempty common intersection.
We shall repeatedly prove compatibility by common refinements or by bounds
on every finite future correction.  An independently existing limit is
never needed as an anchor.

\begin{definition}[Stage schedules]\label{def:stage-schedule}
A stage schedule is an increasing, unbounded map $\sigma:\N\to\N$.
Replacing $X_n$ by $X_{\sigma(n)}$ gives an equivalent computation.
\end{definition}

Thus two algorithms can be compared at different speeds.  A useful named
number can carry several computations connected by proved equivalences;
transitivity supplies comparisons along a finite chain.  There is no reason
to discard a slow geometric construction after discovering a faster series.

\section{Examples of real numbers}

\subsection{Square roots and finite geometry}
\label{rem:algebraic-numbers-sqrt-two}
For a rational $q\ge0$, start with rational $0\le a\le b$ satisfying
$a^2\le q\le b^2$.  Bisect and retain the half whose squared endpoints
still bracket $q$.  If the midpoint square equals $q$, return the exact
constant interval thereafter.  At stage $n$ the width is at most
$(b-a)2^{-n}$, and finite interval multiplication proves that the resulting
computation squares to $q$.

For $\sqrt2$, the rational secant--tangent refinement gives another
construction:
\begin{verbatim}
def sqrt2(n):
    a = 1
    b = 2
    for _ in range(n):
        a = (2 + a*b) / (a + b)
        b = (b + 2/b) / 2
    return [a, b]
\end{verbatim}
All divisions here are exact rational divisions.  For one step, writing
$a',b'$ for the new endpoints,
\[
 (a')^2-2=\frac{(a^2-2)(b^2-2)}{(a+b)^2},\qquad
 (b')^2-2=\frac{(b^2-2)^2}{4b^2}.
\]
Also $a\le a'\le b'\le b$, and
\[
 b'-a'=\frac{(b-a)(b^2-2)}{2b(a+b)}\le\frac{b-a}{2}.
\]
These identities prove validity.  The bisection and secant--tangent
computations are equivalent because both retain the unique nonnegative
solution, with uniqueness following from $(u-v)(u+v)=u^2-v^2$ and the
positive lower bounds in this example.

\begin{figure}[htbp]\centering
\sqrtSecantTangentAnimation
\caption{Secant--tangent rational enclosures for $\sqrt2$.}
\label{fig:sqrt-secant-tangent}
\end{figure}

There is no rational whose square is $2$: in a reduced fraction $p/q$, the
identity $p^2=2q^2$ would make both $p$ and $q$ even.  Thus this is a
nonrational quantity constructed entirely by rational arithmetic.

Square roots also compute lengths of rational vectors.  The finite identity
\[
 (a^2+b^2)(c^2+d^2)-(ac+bd)^2=(ad-bc)^2
\]
proves the coordinate Cauchy--Schwarz inequality and hence the broken-line
inequality.  Shoelace determinants compute polygonal areas exactly.  These
are the geometric tools needed for the circle, rather than a separate
axiomatization of Euclidean space.

\subsection{Other computations of numbers}
\label{rem:sources-of-raw-reals}
The geometric sum illustrates a finite identity with a tail:
\[
 (1-r)\sum_{k=0}^{N}r^k=1-r^{N+1}.
\]
For rational $|r|<1$, the error $|r|^{N+1}/(1-|r|)$ gives a number
computation.  Bounded monotonicity by itself does not supply a shrinking
upper--lower gap; an error estimate or a terminating precision search is
still necessary.

Later we shall construct, and prove equivalent, the computations suggested
by
\[
 \frac\pi4=1-\frac13+\frac15-\cdots
   =\int_0^1\frac{dt}{1+t^2},\qquad
 \sum_{n=1}^{\infty}\frac1{n^2}=\frac{\pi^2}{6}.
\]
At this point these are destinations, not definitions silently imported from
ordinary calculus.  The circle will define $\pi$ geometrically; rectangles,
series, and Fourier sums will then supply different computations of it.

\begin{remark}[Calculus as equivalence]\label{rem:calculus-as-equivalence}
A great deal of calculus is the art of proving that independently constructed
numbers are equivalent.  The fundamental formula
\[
 \int_a^b f'(x)\,dx=f(b)-f(a)
\]
turns an accumulation computation into an endpoint computation.  Its
hypotheses must explain why local derivative information controls a whole
rational interval; pointwise differentiation on the rationals alone does not.
\end{remark}

\section{Complex numbers}

Let $\Q(i)=\{r+is:r,s\in\Q\}$.  A rational complex box is a pair of
rational coordinate intervals.  A complex computation gives nested boxes
whose two widths shrink to zero; equivalence is coordinatewise overlap.
Addition and multiplication are rational box operations, using
\[
 (r+is)(u+iv)=(ru-sv)+i(rv+su).
\]
For a denominator separated from zero use
$z^{-1}=\overline z/( (\Re z)^2+(\Im z)^2)$ and the real reciprocal
construction.  Finite rational coordinate bounds control all errors.

The same convention applies to vectors and matrices: finitely many number
computations, with finite sums and products.  No additional completeness
principle appears when the number of coordinates increases.

\section{Algorithms with a variable}
\label{rem:parameterized-constructions}

\begin{definition}[Function computation]\label{def:real-fun-raw}
A real-valued function computation consists of a declared domain of rational
inputs and an algorithm
\[
 (x,n)\longmapsto [a_f(x,n),b_f(x,n)]
\]
which is a valid number computation at each permitted input $x$.
\end{definition}

The outputs need not be rational.  In particular, $\sin(\pi x)$ will have
rational inputs and computed outputs.  An algebraic, dyadic, or rational
geometric chart may be a more natural first domain than all rationals;
passing between charts requires an explicit comparison of computations.

\begin{definition}[Interval continuity]\label{def:interval-regular}
On a rational interval $[a,b]$, a continuity certificate supplies, for every
positive rational $\varepsilon$, a positive rational $\mu(\varepsilon)$
such that
\[
 |x-y|\le\mu(\varepsilon)\quad\Longrightarrow\quad
 |f(x)-f(y)|\le\varepsilon
\]
for rational $x,y$ in that interval.  This is uniform continuity on the
specified interval, with usable data.  A different interval may have a
different certificate.
\end{definition}

Together with point evaluation, this produces small output enclosures for
small input intervals: evaluate at a rational sample and add the continuity
error on both sides.  Conversely, shrinking interval-image bounds give
this continuity certificate.  A finite grid and coarse value enclosures
also give an absolute bound for $f$ on $[a,b]$.

This is stronger than continuity at each rational point.  We do not infer it
from compactness.  For square roots, finite algebra gives
\[
 |\sqrt{x}-\sqrt{y}|\le\sqrt{|x-y|}\qquad(x,y\ge0),
\]
so $\mu(\varepsilon)=\varepsilon^2$ works.  For a polynomial, expand
$y^k-x^k=(y-x)\sum_{j=0}^{k-1}y^{k-1-j}x^j$ and bound the finite sum on
the interval.  Separated rational functions follow by the reciprocal bound.

\begin{proposition}[Computed inputs]\label{prop:computed-inputs}
An interval-continuous rational-input computation can be evaluated on a
number computation known to stay in its interval.
\end{proposition}
\begin{proof}
Intersect each input enclosure with $[a,b]$, choose a rational sample in
that finite intersection, and request enough input precision that all its
samples have outputs within the desired continuity tolerance.  Evaluate
the sample to the remaining tolerance.  The resulting output boxes are
compatible across input refinements and shrink.  Finite-prefix intersections
make them nested.  Equivalent input computations give equivalent outputs
by the same estimate.
\end{proof}

This is the precise meaning of an implicit continuation beyond rational
inputs.  We have constructed an evaluation procedure, not postulated a
function on a completed line.

\section{Inverting a monotone computation}
\label{thm:constructive-inverse-function}
Suppose $f$ is increasing and interval-continuous on $[a,b]$, and a
positive separation bound $\sigma(d)$ is supplied:
\[
 y-x\ge d>0\quad\Longrightarrow\quad f(y)-f(x)\ge\sigma(d)>0.
\]
Let a computed target $z$ be certified between $f(a)$ and $f(b)$.
Use the two trisection points $u,v$ of the current rational bracket.
Their images are separated by $\sigma(v-u)$.  Refine the value enclosures
until either $z>f(u)$ or $z<f(v)$ is certified.  At least one strict
inequality must become visible, even when the target equals one of the
other values.  Keep $[u,b]$ in the first case and $[a,v]$ in the second.

Each step preserves a range bracket and shrinks the input interval by a
factor $2/3$.  Interval continuity proves that the resulting input
computation maps to $z$.  Separation proves uniqueness.  This is a
terminating inverse construction without a decision procedure for equality
of computed numbers.  Squaring on $[0,1]$ has
$\sigma(d)=d^2$; the angle chart in the next chapter has a linear bound.

\section{Domains and equivalent descriptions}
Functions agree only where both computations are defined.  A geometric
chart, a series disk, and a positive algebraic branch may describe the same
function on overlapping domains without being interchangeable everywhere.
Finite chains of proved comparisons retain these domain restrictions.

The rest of the book uses only the constructions just described: finite
rational operations, enclosures, terminating refinement, and equivalence.
When a new infinite process occurs, its finite future errors will be
bounded explicitly.  No theorem identifying these computations with an
ambient completed number system is required.
